\documentclass[11pt,a4paper]{article}
\usepackage[margin=24mm]{geometry}
\usepackage[T1]{fontenc}
\usepackage{lmodern,amsmath,amssymb,booktabs,longtable,array,microtype}
\usepackage[hidelinks]{hyperref}
\urlstyle{same}
\setlength{\emergencystretch}{2em}
\title{Full Pair-Entropy Study from Existing MD Trajectories\\
\large Protocol for the eight-mixture charged-fluid study}
\author{EntropicFluids collaboration}
\date{31 October 2026}
\begin{document}
\maketitle
\begin{abstract}
This protocol plans an independent MD determination of the complete all-species
pair entropy $s_2$, selected probe-centered contributions, and their relationship
to thermodynamic-integration (TI) entropy and self-diffusion. The primary study
uses the existing eight compositions, three replicas each, at $T^*=2$ and
$\rho^*=0.7$. All-particle positions permit extraction of every partial radial
distribution function (RDF) without new MD. The main challenges are nonlinear
histogram bias, rare ion-pair statistics, finite-size normalization, and radial
convergence. The protocol defines checks and decision criteria before fitting
entropy--diffusion relations. No new production simulations are submitted by
this planning document.
\end{abstract}

\section{Questions, scope, and existing evidence}
The primary question is whether replacing total excess entropy by full pair
entropy preserves the conclusions of our replicated MD study. Secondary
questions are whether the remainder $s_{\rm ex}^{\rm TI}-s_2^{\rm MD}$ changes
with physical solvent diversification, and whether ion-centered structural
contributions explain probe mobility better than the solvent-dominated bulk
sum. Rosenfeld-type scaling~\cite{Rosenfeld1977} and its mixture/tracer
extensions~\cite{Krekelberg2009} motivate these tests, but no universal relation
is assumed. The local pair statistic described by He et al.~\cite{He2026}
motivates a separate selected-pair comparison; it is not interchangeable
with the full mixture entropy.

Use all eight states B3, B5, U4, L5, M5, H5, U6, and H7, with three replicas
per state. All contain 120 particles of each divalent sign, out of $N=4000$.
Exact compositions and type definitions are in
\path{ComponentEntropyStudy_20261030/selected.json}; the composition tables
are reproduced in the existing reports. The primary contrasts are L5--U4
(label null), M5--U4 and H5--U4 (moderate and strong solvent diversification),
H7--U6 (diversification with monovalent background), and B5--B3 (salt addition).
Broad baseline-to-design comparisons are secondary.

Existing HNC outputs preserve the same qualitative ordering for full pair
and total entropy. For H5--U4 and H7--U6, HNC pair increments are about
$0.075k_B$ versus total HNC increments near $0.111k_B$. These preliminary
numbers are recorded in \path{doc/hnc_pair_entropy_comparison.json}.
They are hypotheses for an MD test, not MD pair-entropy measurements.
The completed diffusion estimates, including the provisional approximately
5\% controlled probe response, remain unchanged by this analysis.

\section{Input data and feasibility}
Each dense production trajectory is located under the study directory at
\begin{quote}\small
\path{md/<case>/replica_<r>/transport/positions.nc}.
\end{quote}
The files contain all-particle \path{unwrapped_coordinates}, IDs, types,
frame times, cell origins, lengths, and angles. A header inventory is stored
in \path{doc/md_pair_entropy_trajectory_inventory.json}. The standard record
has 2001 frames, $N=4000$, and positions every $0.5\tau$ over $1000\tau$.
The cubic box length is approximately $17.878\sigma$, giving a maximum
complete spherical radius of approximately $8.939\sigma$.

Use these position records, not the velocity files: the latter contain only
subsets of neutral species. Resolve types and charges against the immutable
input \path{transport/data.mixture}, \path{selected.json}, and replica
manifest. Check frame IDs and sorting explicitly. Read time scale factors
and box metadata through the NetCDF interface rather than assuming raw
array values are already physical times.

Primary structural averages use the final $500\tau$ of dense production,
matching the TI averaging interval at the dense endpoint. Analyze the full
$1000\tau$ and its two halves as stationarity/sampling sensitivities. First
benchmark 20 frames from one replica of B3 and H7; then process all replicas
with matched frame counts. Increase from 50 to 100, 200, 500, and, where
needed, all approximately 1001 frames in the final half. All particles are
retained; only frame thinning is allowed in the initial sampling pilot.

A direct $N=4000$ calculation has 7,998,000 unordered pairs per frame.
Full-frame analysis across 24 replicas therefore requires substantial
postprocessing despite requiring no new dynamics. Measure time, memory,
and I/O in the pilot before selecting CPU allocation or an offline GPU
implementation. Stream coordinates and accumulate histograms in a compiled
pair-distance kernel; avoid storing all distances or loading every trajectory
at once. Parallelize independent replicas within measured resource limits.
Do not rerun forces or electrostatics merely to obtain distance histograms.

\section{RDF estimation and finite-system conventions}
Wrap positions into the periodic box and calculate minimum-image separations.
This is a static observable, so no COM subtraction is needed. Only the
subsequent transport comparison uses the established common-system COM frame.
Exclude self-pairs and count each physical pair once.

For a bin $[r_k,r_{k+1})$, use its exact spherical shell volume
\begin{equation}
\Delta V_k=\frac{4\pi}{3}(r_{k+1}^3-r_k^3).
\end{equation}
If $H_{ab,k}$ is the unordered histogram accumulated over $F$ frames of a
fixed-volume system, define the ideal-normalized finite-system RDF as
\begin{equation}
g_{ab,k}=\frac{H_{ab,k}}{F M_{ab}\Delta V_k/V},\qquad
M_{ab}=\begin{cases}N_aN_b,&a\ne b,\\
N_a(N_a-1)/2,&a=b.\end{cases}
\end{equation}
This convention gives unity in expectation for ideal random positions at
fixed species counts. Document any conversion to the alternative bulk-density
normalization, especially for small ionic populations. Do not let a diagonal
$1-1/N_a$ offset accumulate as an apparent long-range entropy contribution.
If the box varies, accumulate the frame-dependent ideal counts explicitly;
our intended records are constant-volume and must be checked as such.

Compute bins from zero to a common radius below the smallest half-box length.
Use $\Delta r=0.02\sigma$ initially, with $0.01$, $0.04$, and where useful
$0.08\sigma$ as resolution checks. Do not extend the spherical normalization
beyond half the box. Record both raw counts and normalized RDFs, including
metadata and counts for every species pair. Do not average RDFs from different
replicas before estimating replica uncertainty.

Mechanically identical secondary solvents in L5 must be merged for the primary
physical-species calculation. Independently process an explicitly label-resolved
version as a diagnostic. Nonlinear entropy estimates from noisy, subdivided
histograms need not be exactly label invariant even when interactions are;
merging the counts before estimating the physical RDF avoids that artifact.
Check actual U4/L5 positional identity by ID, and also perform an exact
synthetic split/merge on the same snapshots.

\section{Entropy integrals and pair decomposition}
For $I_{ab}(R)=\int_0^R\phi(g_{ab}(r))r^2\,dr$, where
$\phi(g)=g\ln g-g+1$, calculate the full pair entropy
\begin{equation}
\frac{s_2(R)}{k_B}=-2\pi\rho\sum_{ab}x_ax_b I_{ab}(R).
\end{equation}
The sum is over ordered species pairs. With an unordered pair list, store
\begin{equation}
C_{aa}=-2\pi\rho x_a^2 I_{aa},\qquad
C_{ab}=-4\pi\rho x_ax_b I_{ab}\quad(a<b),
\end{equation}
so that $s_2/k_B=\sum_{a\le b}C_{ab}$. This makes the cross-pair factor of two
explicit. For histogram RDFs, integrate using exact radial bin weights
$(r_{k+1}^3-r_k^3)/3$ rather than an inconsistent midpoint volume.
At $g=0$, use $g\ln g=0$ and $\phi(0)=1$; do not discard empty core bins or
apply an arbitrary density floor.

Store full cumulative integrals and decompose them into neutral--neutral,
neutral--ion, ion--ion, divalent-containing, and monovalent-containing pairs.
Categories must be defined without overlap when their sum is compared with
the total. A useful additional species-centered allocation is
\begin{equation}
\frac{s_2^{(a)}}{k_B}=-2\pi\rho\sum_bx_b I_{ab},\qquad
s_2=\sum_ax_as_2^{(a)}.
\end{equation}
These are structural allocations, not partial molar entropy derivatives.
Report $+2$ and $-2$ separately and their equal-population average.

Define selected probe metrics carefully. The weighted central-probe allocation
$\sum_{a\in\{+2,-2\}}x_as_2^{(a)}$ counts a probe--neutral term once in
the ordered sum, whereas the union of all unordered pairs involving a probe
contains both ordered directions of that cross term. They are distinct
observables. Report both with explicit definitions rather than calling both
local excess entropy. Compare them to the full sum and to the He et al.
selected-pair construction only after resolving its normalization.

\section{Bias, radial convergence, and mandatory validation}
The function $\phi(g)$ is convex. Histogram noise therefore creates a negative
bias in $s_2$, particularly in sparsely populated ionic bins and nominally flat
tails. A small statistical error bar does not eliminate this systematic bias.
Use a pooled RDF within each replica for the primary estimate: the entropy of
the pooled RDF and the mean of entropies of small-block RDFs are different
estimators. Blocks diagnose uncertainty and sample-size bias; the latter mean
is not substituted silently for the former.

Required checks are:
\begin{enumerate}
\item Independent ideal random-point samples, with the same box, counts,
frame counts, binning, and normalization. The expectation of the RDF is one,
but the plug-in entropy is generally negatively biased at finite sample size.
Quantify this bias instead of assuming the observed ideal-gas entropy is zero.
Do not subtract its measured bias from interacting data without validation.
\item A small-frame independent direct-distance/histogram cross-check;
translation, atom-order and type-code invariance; RDF symmetry;
self-pair exclusion; exact pair-count bookkeeping; and recovery of the
single-species formula when physical species are merged.
\item Nonpositivity of every converged pair contribution, correct hard-core
integration, total/species allocation additivity, and the exact synthetic
label split/merge null. Statistical differences of label-resolved plug-in
estimates must not be interpreted as physical entropy changes.
\item Sample-count and bin-width convergence, raw rare-bin counts, first/second
half stationarity, and effective correlation times. More closely spaced frames
are not automatically independent configurations.
\item Cumulative entropy curves $s_2(R)$ and pair-resolved tails. Compare
$R=4,5,6,7,8\sigma$ and the allowed maximum, with particular attention to ionic
pairs and the small controlled increments. Do not set $g=1$ beyond a chosen
radius merely to force a plateau.
\end{enumerate}

An optional smooth or fitted tail must be fixed independently of TI or diffusion,
validated on held-out radii, and reported separately from the directly sampled
integral. HNC can guide a tail sensitivity but cannot serve as independent
MD confirmation. If no radial plateau is resolved within the box, report a
truncated $s_2(R)$ with a systematic limitation and plan a larger box.
Also compare plausible finite-system RDF normalization conventions; an
ideal-normalized RDF does not remove all finite-size effects in correlations.

\section{Statistical analysis and comparison with TI and diffusion}
Produce one converged estimate per independent replica. Use paired composition
contrasts across the three matched replica indices, with nominal 95\%
Student-$t$ intervals and two degrees of freedom, consistent with the current
study. Within-replica block resampling or a validated block jackknife can
quantify histogram uncertainty and bias, but cannot create additional MD
replicas. Correlated frame or pair counts are not the sample size for the
between-replica confidence interval. Retain covariance when comparing TI and
$s_2$ from the same replica.

The main thermodynamic outputs are $s_2^{\rm MD}$, $s_{\rm ex}^{\rm TI}$,
$r_{\ge3}=s_{\rm ex}^{\rm TI}-s_2^{\rm MD}$, and their paired increments.
Interpret $r_{\ge3}$ as the residual beyond the pair approximation under
matched conventions; unresolved RDF, TI, finite-size, and tail errors also
enter it. Evaluate $s_2/s_{\rm ex}$ only with uncertainty and a common
normalization. Do not impose a 90\% pair fraction. Compare HNC and MD RDFs
and $s_2$ at the same density, temperature, charges, composition, and potential
convention; verify the input LJ shift/cutoff and numerical radial conventions.
Do not mix HNC $s_2$ with MD total entropy when estimating the MD residual.

For dynamics retain the existing primary 20--100$\tau$ MSD fitting window,
with 10--50 and 50--200$\tau$ checks. Compare $\ln D^{\rm R}$ with full TI
entropy, full MD $s_2$, and the predefined probe metrics. Use
$D_i^{\rm R}=D_i\rho^{1/3}\sqrt{m_i/(k_BT)}$, with total number density.
Assess mean self-diffusion $\sum_ix_iD_i$, each divalent sign, and the pooled
probe mean separately. These are not conductivity or mutual diffusion.

Exclude L5 from regression counts; fit fixed-ionic-strength composition families
separately. With only seven physical states, restrict initial models to one
entropy coordinate plus an intercept, show residuals and paired contrasts,
and use leave-one-state-out or cross-background prediction. Uncertainty in
both entropy and diffusion and sensitivity to individual states must be shown.
A high $R^2$ from a few composition means is not evidence of a universal law.
No RDF smoothing, radius, or structural definition is selected to maximize
agreement with diffusion.

\section{Execution stages, decisions, and deliverables}
\begin{longtable}{@{}p{0.23\textwidth}p{0.71\textwidth}@{}}
\toprule
Stage & Work and decision\\
\midrule
\endhead
A: inventory/pilot & Verify all 24 files, IDs, physical species, time windows,
and boxes. Benchmark B3 and H7, establish a tested histogram kernel, and pass
ideal-gas and synthetic-label checks before full processing.\\[3pt]
B: all-pair RDFs & Process all eight states and three replicas. Increase frame
counts and vary bins/radii until the selected contrasts are stable or a
sampling limitation is demonstrated. Preserve histograms for reanalysis.\\[3pt]
C: entropy benchmark & Compare MD $s_2$ and selected contributions with TI and
HNC. Proposed practical targets are a 95\% half-width below $0.005k_B$ for
controlled entropy increments and bin/frame/radial sensitivity below
$0.002k_B$. These are design goals, not guarantees or a formal equivalence test.\\[3pt]
D: transport assessment & Test whether entropy ordering, controlled contrasts,
and descriptive diffusion relationships survive replacement of TI by pair
entropy. Changing an entropy estimator does not resolve the existing uncertain
5\% probe diffusion response.\\[3pt]
E: conditional new MD & If sampling is limiting, extend dense records or add
replicas. If radial/finite-size errors dominate, compare a second size, for
example $N=8000$ at the same state and fractions, using the same physical LJ
potential and controlled electrostatic accuracy. Submit only after a separate
execution decision; no new jobs are required for the initial extraction.\\
\bottomrule
\end{longtable}

The dense trajectories suffice for the primary entropy--diffusion question.
Positions at all 19 TI densities are not established as available: the
production writer records all-particle position trajectories at the dense
transport stage. A density-dependent $s_2(\rho)$ extension therefore requires
an inventory of any additional stored configurations or new sampling.
Prioritize a few representative densities for B3, U4, H5, U6, and H7 only if
the dense benchmark indicates a scientific need. Higher concentration and
size/charge asymmetry are later transfer studies, not prerequisites for the
current all-pair calculation.

Write a separate analysis directory, for example
\path{MDPairEntropyAssessment/}, containing a protocol/configuration file,
source and trajectory checksums, software versions, per-replica raw pair
counts and RDFs, cumulative pair integrals, full and selected $s_2$ values,
noise/tail/bin/frame checks, paired contrast intervals, TI residuals, and
transport-comparison plots. Preserve machine-readable CSV/JSON summaries
and the exact frame manifest. The final report should explicitly decide
whether pair entropy changes the qualitative conclusions, quantitative
increments, or species-versus-bulk interpretation, and identify any questions
that remain unresolved rather than forcing agreement with the HNC expectation.

\begin{thebibliography}{9}
\small
\bibitem{Rosenfeld1977}
Y. Rosenfeld, \textit{Phys. Rev. A} \textbf{15}, 2545--2549 (1977).
\href{https://doi.org/10.1103/PhysRevA.15.2545}{doi:10.1103/PhysRevA.15.2545}.
\bibitem{Krekelberg2009}
W. P. Krekelberg et al., ``Generalized Rosenfeld scalings for tracer diffusivities
in not-so-simple fluids: Mixtures and soft particles,'' \textit{Phys. Rev. E}
\textbf{80}, 061205 (2009).
\href{https://doi.org/10.1103/PhysRevE.80.061205}{doi:10.1103/PhysRevE.80.061205}.
\bibitem{He2026}
L. He et al., ``Bridging thermodynamics and ion transport in lithium batteries
via molecular design of high-entropy electrolytes,''
\textit{Molecular Systems Design \& Engineering} \textbf{11}, 681--690 (2026).
\href{https://doi.org/10.1039/D6ME00060F}{doi:10.1039/D6ME00060F}.
See the full-text assessment in \path{doc/entropy_diffusion_prior_work.tex}.
\end{thebibliography}
\end{document}
